\documentclass[11pt]{article}
\usepackage[margin=0.82in]{geometry}
\usepackage{amsmath,amssymb,mathtools,booktabs,array,hyperref,microtype,enumitem}
\usepackage[T1]{fontenc}
\usepackage{lmodern}
\hypersetup{colorlinks=true,linkcolor=blue,citecolor=blue,urlcolor=blue}
\setlength{\parindent}{0pt}
\setlength{\parskip}{0.46em}
\setlist{nosep,leftmargin=*}
\newcommand{\Dx}{D_x}
\newcommand{\Dy}{D_y}
\title{Exact Reduction Certificates for Laurent Periods\\[1mm]
\large The $G,U,V,J,A,\Xi$ formalism}
\author{Bradley Klee and Mech.An.ika Sol (Open AI)}
\date{August 5, 2026\\Public draft with one worked example}

\begin{document}
\maketitle

\begin{abstract}
We consider Laurent polynomials
$F\in\mathbb Q[x^{\pm1},y^{\pm1}]$ and their constant-term power series
\[
  \Phi(t)
  =
  \left[\frac{1}{1-tF}\right]_0
  =
  \sum_{n\ge0}[F^n]_0t^n.
\]
A finite exact reduction constructs a differential annihilator $A$ for
$\Phi$ together with an explicit certificate $\Xi=(\Xi_x,\Xi_y)$.
The general symbolic implementation uses the $G,U,V,J$ formalism developed in
Chapter~3 of Bradley Klee's dissertation.  The coefficient matrix $G$
records the reduction identity, its solution gives the reduced part $U$
and certificate data $V$, and $J$ records the derivative correction made
when the pole order is lowered.  Linear dependence among the resulting
finite-dimensional data gives $A$, while the accumulated $V$-data give
$\Xi$ without a second solve.  The worked public example is the cubic
Hamiltonian recorded with OEIS A303790; its factorization gives a shorter
Laurent telescoper, paired with an independent Hamiltonian certificate.  The
accompanying exact data reproduce its coefficient sequence, differential
equation, recurrence, and both certificate checks.
\end{abstract}

\section{Formal reductions}

Put $K=\mathbb Q(t)$ and
\[
  \rho=1-tF,
  \qquad
  \Dx=x\frac{\partial}{\partial x},
  \qquad
  \Dy=y\frac{\partial}{\partial y}.
\]
The coefficient of $x^0y^0$ in either $\Dx(R)$ or $\Dy(R)$ is zero.
Thus a pair of rational functions $\Xi=(\Xi_x,\Xi_y)$ satisfying
\begin{equation}\label{eq:certificate-goal}
  A\!\left(\frac1\rho\right)=\Dx(\Xi_x)+\Dy(\Xi_y)
\end{equation}
proves $A\Phi=0$.

Fix a positive integer $r$.  Let $w,a,b_x,b_y$ be Laurent
polynomials.  If
\begin{equation}\label{eq:scalar-split}
  w=\rho a-\Dx(\rho)b_x-\Dy(\rho)b_y,
\end{equation}
then the product rule gives the reduction identity
\begin{align}
  \frac{w}{\rho^{r+1}}
  &={}
  \frac{a-\frac1r\bigl(\Dx(b_x)+\Dy(b_y)\bigr)}{\rho^r}
  +\Dx\!\left(\frac{b_x}{r\rho^r}\right)
  +\Dy\!\left(\frac{b_y}{r\rho^r}\right).
  \label{eq:scalar-lowering}
\end{align}
Thus lowering the pole order of $w/\rho^{r+1}$ reduces to solving
\eqref{eq:scalar-split} for $(a,b_x,b_y)$.  The relevant linear map is
\begin{equation}\label{eq:coefficient-map}
  (a,b_x,b_y)\longmapsto
  \rho a-\Dx(\rho)b_x-\Dy(\rho)b_y.
\end{equation}

Let $B=(b_1,\ldots,b_d)$ be the ordered
Laurent-monomial basis for each of $a,b_x,b_y$.  Let
$C=(c_1,\ldots,c_m)$ be an ordered Laurent-monomial basis large enough
to contain $B$, $FB$, $\Dx(F)B$, and $\Dy(F)B$.  For a coefficient
column $q=(q_1,\ldots,q_d)^T\in K^d$, write
\[
  \langle q\rangle=\sum_{j=1}^d q_jb_j.
\]
Angle brackets always reconstruct a Laurent polynomial in the basis $B$.
The rows of $G$ are indexed by the basis $C$, and its columns are indexed
by three ordered copies of $B$, corresponding to $a,b_x,b_y$.

The matrix
\[
  E\in K^{m\times d}
\]
is the inclusion of the $B$-space into the $C$-space.  Its $j$th column
is $b_j$ written in the basis $C$.

The matrix
\[
  G\in K^{m\times3d}
\]
represents the map \eqref{eq:coefficient-map}.  We compute
\[
  U,V_x,V_y\in K^{d\times d}
\]
from the exact matrix identity
\begin{equation}\label{eq:G-inversion}
  G
  \begin{pmatrix}U\\V_x\\V_y\end{pmatrix}=E.
\end{equation}
The stacked matrix has size $3d\times d$, so both sides of
\eqref{eq:G-inversion} have size $m\times d$.

For any $q\in K^d$, multiplication of \eqref{eq:G-inversion} by $q$
gives
\[
  G
  \begin{pmatrix}Uq\\V_xq\\V_yq\end{pmatrix}=Eq.
\]
Therefore the Laurent polynomial $w=\langle q\rangle$ is split by
\eqref{eq:scalar-split} with
\[
  a=\langle Uq\rangle,
  \qquad
  b_x=\langle V_xq\rangle,
  \qquad
  b_y=\langle V_yq\rangle.
\]

If $G$ is square and nonsingular, then
\[
  \begin{pmatrix}U\\V_x\\V_y\end{pmatrix}=G^{-1}E.
\]
If $G$ is rectangular, exact Gauss--Jordan elimination is applied to the
augmented system.  Free variables are set to zero, and the resulting
solution is accepted only when substitution gives zero residual in
\eqref{eq:G-inversion}.

Put
\[
  V=\begin{pmatrix}V_x\\V_y\end{pmatrix}\in K^{2d\times d}.
\]
The matrix
\[
  J\in K^{d\times2d}
\]
represents differentiation of the certificate numerators:
\[
  \left\langle J
  \begin{pmatrix}q_x\\q_y\end{pmatrix}\right\rangle
  =\Dx\langle q_x\rangle+\Dy\langle q_y\rangle.
\]
Substitution in \eqref{eq:scalar-lowering} gives the one-step reduction
matrix
\begin{equation}\label{eq:Lower}
  \operatorname{Lower}_r=U-\frac1rJV\in K^{d\times d}.
\end{equation}
A kernel relation among reduced coefficient columns gives $A$; the same
linear combination of the stored $V_x,V_y$ columns gives $\Xi$.

\section{Example}

\textbf{The cubic Hamiltonian and OEIS A303790.}

Let
\begin{equation}\label{eq:A303790-K}
  K(p,q)=p^2+q^2+p^3+q^3,
  \qquad t=E/32.
\end{equation}
The normalized period at either finite center has the form
\[
  \Pi(E)=\sum_{n\ge0}A_n(E/32)^n.
\]
The equality of the two local coefficient sequences follows from the affine
involution $(p,q)\mapsto(-p-2/3,-q-2/3)$, which sends $K$ to $8/27-K$.

Define the integral Laurent polynomial
\begin{equation}\label{eq:A303790-F}
  F_{303790}(w,y)=
  \frac{(w+1)^3(y+1)^2(y^2-4y+1)^2}{w^2y^3}.
\end{equation}
Lagrange inversion for the radial equation gives the exact coefficient
identity
\begin{equation}\label{eq:A303790-CT}
  A_n=[F_{303790}(w,y)^n]_0.
\end{equation}
Thus integrality is immediate.

For the exact recurrence proof, factor $F_{303790}=B(w)C(y)$.  The
$w$-factor gives $[B^n]_0=\binom{3n}{2n}$.  For the $y$-factor, the supplied
rational function $R(n,y)$ satisfies
\begin{equation}\label{eq:A303790-R}
  P_0(n)+P_1(n)C+P_2(n)C^2
  =y\frac{\partial R}{\partial y}
   +n\,y\frac{C'}{C}R.
\end{equation}
After multiplication by $C^n$, the right side is
$y\,\partial_y(RC^n)$ and has zero constant term.  Combining the resulting
recurrence for $[C^n]_0$ with the exact ratio for $\binom{3n}{2n}$ gives
\begin{equation}\label{eq:A303790-recurrence}
  (n+1)^2A_{n+1}
  -12(27n^2+27n+5)A_n
  +2592(3n-2)(3n+2)A_{n-1}=0.
\end{equation}
The full $P_0,P_1,P_2,R$ data are kept in the accompanying exact certificate.

For $Y(t)=\sum_{n\ge0}A_nt^n$, this is equivalent to
\begin{equation}\label{eq:A303790-ODE}
\begin{aligned}
  &t(108t-1)(216t-1)Y''
  +(69984t^2-648t+1)Y'\\
  &\hspace{43mm}+60(216t-1)Y=0.
\end{aligned}
\end{equation}
The independent Hamiltonian certificate proves
\[
  E(27E-8)(27E-4)\Pi''
  +(2187E^2-648E+32)\Pi'
  +15(27E-4)\Pi=0.
\]
Substitution of $E=32t$ gives exactly \eqref{eq:A303790-ODE}.  Both certificate
residuals are zero by exact coefficient reduction.  The first values are
\[
  1,60,7380,1090320,176978340,30471320880,\ldots,
\]
in agreement with OEIS A303790.

\clearpage
\section{General $G,U,V,J$ pseudocode}

The public A303790 packet also contains the shorter factorized replay used in
the worked example.  The following is the general matrix-reduction route.

\begingroup
\fontsize{9.5}{10.9}\selectfont
\begin{enumerate}[label=\texttt{\arabic*.},leftmargin=2.7em,itemsep=0.22em,topsep=0pt]
\item Read $F$ as an exact Laurent polynomial; set $\rho=1-tF$ and list the exponent pairs occurring in $1$ and $F$.
\item Form every pairwise sum of those exponent pairs; let $B$ be all integer points in their convex hull, ordered lexicographically.
\item Let $C$ be the union of $B$ with every translate $B+(i,j)$ for which $x^iy^j$ occurs in $F$; order $C$ lexicographically.
\item Represent Laurent polynomials on $B$ or $C$ by exact coefficient dictionaries keyed by exponent pairs.
\item For each $b_j$ in $B$, place the $C$-coordinates of $\rho b_j$, $-\Dx(\rho)b_j$, and $-\Dy(\rho)b_j$ in columns $j$, $d+j$, and $2d+j$ of $G\in K^{m\times3d}$.
\item Build $E\in K^{m\times d}$ and $J\in K^{d\times2d}$ in the ordered bases $B,C$.
\item Solve $G(U,V_x,V_y)^T=E$ by exact Gauss--Jordan elimination, or solve only $G(u,v_x,v_y)^T=Ef_2$ when only the $F^2$ column is required; then impose $\operatorname{Lower}_2f_2=\alpha r+\beta f$.
\item Evaluate this system at the first regular integer value of $t$; perform exact Gauss--Jordan elimination.
\item At successive regular integer values of $t$, solve that same square block exactly over $\mathbb Q$; keep training values and three held-out values.
\item For each unknown coefficient, reconstruct the rational function of least total numerator-plus-denominator degree fitting the training values; reject it unless it also fits every held-out value.
\item Set all free unknowns to zero, rebuild the symbolic solution, and substitute it into the original block system over $\mathbb Q(t)$.
\item Clear stated denominators, expand in the declared monomial basis, and require every coefficient of every residual to be zero.
\item Extract the selected column vectors $u,v_x,v_y$ and $\alpha,\beta$; verify $G(u,v_x,v_y)^T=Ef_2$ and $\operatorname{Lower}_2f_2=\alpha r+\beta f$.
\item Form $X=\bigl(\begin{smallmatrix}1&0&2\alpha\\0&1&2\beta\end{smallmatrix}\bigr)$; compute a nonzero kernel vector and normalize it by denominator clearing, polynomial-gcd removal, integer-content removal, and sign fixing.
\item Write the normalized vector as $(p_0,p_1,p_2)^T$; set $A=p_0+p_1\partial_t+p_2\partial_t^2$ and $\Xi=(p_2\langle v_x\rangle/\rho^2,p_2\langle v_y\rangle/\rho^2)$.
\item Verify $Xp=0$ and verify the cleared coefficient identity corresponding to
      $A(1/\rho)=\Dx(\Xi_x)+\Dy(\Xi_y)$.
\item Convert $A$ to $\sum_k t^kQ_k(\theta)$ with $\theta=t\partial_t$; output
      $\sum_kQ_k(n-k)[F^{n-k}]_0=0$.
\item Generate the requested constant terms by lattice convolution; test the recurrence on all available indices and compare the first five values with direct power expansion.
\item Serialize the ordered bases, sparse matrices, reduction data, $A$, $\Xi$, recurrence, constant terms, and exact pass/fail checks; verify all declared dimensions and nonzero counts.
\item Report elapsed time after basis construction, matrix construction, the exact solve, reduction replay, certificate replay, term generation, and serialization; stop at the first failed check.
\item Do not call a generic simplification routine.  Use only named operations whose effects are visible in the proof record.
\end{enumerate}
\endgroup

\clearpage
\section{Algorithm and data model}

The public package contains the exact A303790 replay, the eleven-record
public catalogue, the catalogue verifier, the generic $G,U,V,J$ implementation,
and its JSON Schema.  The Python source and schema remain separate files rather
than typeset appendices.  This PDF explains the calculation; the package files
supply the literal executable identities and coefficient data.  No private
reference record is copied into the public artifact.

The public replay verifies every stored coefficient list and differential
equation, then repeats both exact certificates for each completed catalogue
entry.  The A303790 scalar and Laurent certificates are also replayed by their
independent scripts.  The generic $G,U,V,J$ verifier takes an explicit data root,
rebuilds $G$ and $J$ from the stated Laurent polynomial, and writes its report
beside the selected data rather than into the code directory.

A complete $G,U,V,J$ certificate record contains the following information:
\begin{center}
\small
\begin{tabular}{>{\ttfamily}p{0.28\textwidth}p{0.62\textwidth}}
\toprule
schema & format name and version\\
model & stable model name\\
input & $F$, $\rho$, coefficient field, basis rule, and requested term counts\\
basis & the ordered exponent pairs used for the unknown and output spaces\\
matrices & sparse exact records for $G$ and $J$\\
reduction & matrix dimensions, selected column vectors $u,v_x,v_y$, $\alpha,\beta$, and the reduced column\\
reduced\_columns\_X & the two-row matrix whose kernel gives $A$\\
annihilator\_A & derivative form and $\theta$ form\\
certificate\_Xi & complete rational $\Xi_x,\Xi_y$, stored numerators, denominator power, and the exact identity\\
recurrence & the polynomials $Q_k$ and the indexing convention\\
constant\_terms & independently generated reference values\\
checks & exact replay results, all required to be true\\
\bottomrule
\end{tabular}
\end{center}
Matrices and vectors are stored with their dimensions, deterministic labels,
declared nonzero counts, and sorted sparse coordinates.  Rational expressions
are stored after explicit denominator collection and cancellation.  Unknown
fields are rejected, and a publication record is accepted only when every
required proof check is true.

The principal public commands, run from the project root, are
\begin{verbatim}
python code/run_examples.py public
python code/build_public_release.py Laurent_period_public_release.zip
\end{verbatim}
The first command writes its exact report below \texttt{examples/public/}; the
second uses an explicit allowlist and rejects every private data path.  Timing
measurements are kept in the root README rather than in the mathematical
certificate records.

\begin{thebibliography}{9}

\bibitem{Klee}
B. J. Klee,
\emph{An Update on the Computational Theory of Hamiltonian Period Functions},
Ph.D. dissertation, University of Arkansas, Fayetteville, December 2020,
especially Chapter~3.

\bibitem{BDS}
A. Bostan, L. Dumont, and B. Salvy,
``Algebraic Diagonals and Walks: Algorithms, Bounds, Complexity,''
\emph{Journal of Symbolic Computation} 83 (2017), 68--92.

\bibitem{BCChL}
A. Bostan, S. Chen, F. Chyzak, and Z. Li,
``Complexity of creative telescoping for bivariate rational functions,''
\emph{Proceedings of ISSAC 2010}, 203--210.

\bibitem{AZ}
G. Almkvist and D. Zeilberger,
``The method of differentiating under the integral sign,''
\emph{Journal of Symbolic Computation} 10 (1990), 571--591.

\bibitem{A303790}
OEIS Foundation Inc., \emph{The On-Line Encyclopedia of Integer Sequences},
A303790.

\bibitem{SymPy}
A. Meurer et al.,
``SymPy: symbolic computing in Python,''
\emph{PeerJ Computer Science} 3 (2017), e103.

\end{thebibliography}

\end{document}
